% part: intuitionistic-logic
% chapter: soundness-completeness
% section: canonical-model

\documentclass[../../../include/open-logic-section]{subfiles}

\begin{document}

\olfileid{int}{sc}{mod}

\olsection{النموذج القانوني}

نختار لعوالم هذا النموذج متتاليات منتهية~$\sigma$ من الأعداد الطبيعية،
بحيث~$\sigma \in \Nat^*$. وتعريف~$\Nat^*$ استقرائي على الوجه الآتي:
\begin{enumerate}
\item $\emptyseq \in \Nat^*$.
\item متى كان~$\sigma \in \Nat^*$ و$\OLId{latin-ordinary}{n} \in \Nat$، كان
  $\sigma.\OLId{latin-ordinary}{n} \in \Nat^*$؛ ونعني بـ~$\sigma.\OLId{latin-ordinary}{n}$
  $\sigma \concat \tuple{\OLId{latin-ordinary}{n}}$، وبـ~$\sigma \concat \sigma'$ وصل
  $\sigma$ بـ~$\sigma'$.
\item وليس في~$\Nat^*$ شيء وراء ما يقتضيه هذان الشرطان.
\end{enumerate}
ولهذا جاز أن تكون~$\Nat^*$ أساسًا لتعريفات استقرائية.

ولتكن~$\tuple{!B_\OLMathNumeral{1}, !C_\OLMathNumeral{1}}$، $\tuple{!B_\OLMathNumeral{2}, !C_\OLMathNumeral{2}}$، \dots، تعدادًا
لأزواج ال!!{formula}s كلها. إذا أُعطيت مجموعة
ال!!{formula}s~$\OLId{greek-variable}{Delta}$، عرّفنا~$\OLId{greek-variable}{Delta}(\sigma)$ بالاستقراء التالي:
\begin{enumerate}
\item $\OLId{greek-variable}{Delta}(\emptyseq) = \OLId{greek-variable}{Delta}$.
\item $\OLId{greek-variable}{Delta}(\sigma.\OLId{latin-ordinary}{n}) = {}$
  \[
  \begin{cases}
    (\OLId{greek-variable}{Delta}(\sigma) \cup \{!B_\OLId{latin-ordinary}{n}\})^* &
    \text{إذا كان $\Delta(\sigma) \cup \{!B_n\} \Proves/ !C_n$} \\
    \OLId{greek-variable}{Delta}(\sigma) & \text{خلاف ذلك}
  \end{cases}
  \]
\end{enumerate}
والمراد بـ~$(\OLId{greek-variable}{Delta}(\sigma) \cup \{!B_\OLId{latin-ordinary}{n}\})^*$ مجموعة
ال!!{formula}s الأولية التي نحصل عليها من
\olref[lin]{lem:lindenbaum} حين نطبقها على
$\OLId{greek-variable}{Delta}(\sigma) \cup \{!B_\OLId{latin-ordinary}{n}\}$ وعلى ال!!{formula}~$!C_\OLId{latin-ordinary}{n}$. ومن مقتضى
هذا الاختيار أنه متى كان~$\OLId{greek-variable}{Delta}(\sigma) \cup \{!B_\OLId{latin-ordinary}{n}\} \Proves/ !C_\OLId{latin-ordinary}{n}$
ثبت~$\OLId{greek-variable}{Delta}(\sigma.\OLId{latin-ordinary}{n}) \Proves !B_\OLId{latin-ordinary}{n}$ وبقي
$\OLId{greek-variable}{Delta}(\sigma.\OLId{latin-ordinary}{n}) \Proves/ !C_\OLId{latin-ordinary}{n}$. كذلك يتحقق
$\OLId{greek-variable}{Delta}(\sigma) \subseteq \OLId{greek-variable}{Delta}(\sigma.\OLId{latin-ordinary}{n})$ لكل~$\OLId{latin-ordinary}{n}$. فإذا كانت
$\OLId{greek-variable}{Delta}$ أولية، لزم أن تكون~$\OLId{greek-variable}{Delta}(\sigma)$ أولية لكل~$\sigma$.

\begin{defn}\ollabel{defn:canonical-model}
  لتكن~$\OLId{greek-variable}{Delta}$ أولية. نسمّي \emph{النموذج القانوني}
  $\mModel{M(\Delta)}$ للمجموعة~$\OLId{greek-variable}{Delta}$ النموذج الذي أجزاؤه كما يأتي:
  \begin{enumerate}
  \item $\OLId{latin-looped}{W} = \Nat^*$، أي مجموعة المتتاليات المنتهية من الأعداد الطبيعية.
  \item $\OLId{latin-looped}{R}$ ترتيب جزئي يحدّه التكافؤ الآتي:~$\OLId{latin-looped}{R}\sigma\sigma'$ إذا وفقط إذا
    كانت~$\sigma$ قطعة ابتدائية من~$\sigma'$؛ ومعنى ذلك أن
    $\sigma'=\sigma\concat\sigma''$ لمتتالية~$\sigma''$ ما.
  \item $\OLId{latin-looped}{V}(\OLId{latin-ordinary}{p}) = \Setabs{\sigma}{\OLId{latin-ordinary}{p} \in \OLId{greek-variable}{Delta}(\sigma)}$.
  \end{enumerate}
\end{defn}

والتحقق من كون~$\OLId{latin-looped}{R}$ ترتيبًا جزئيًا يسير. وأما شرط الرتابة في~$\OLId{latin-looped}{V}$
فحاصل أيضًا: إذ من~$\OLId{greek-variable}{Delta}(\sigma)\subseteq\OLId{greek-variable}{Delta}(\sigma.\OLId{latin-ordinary}{n})$ يلزم بالاستقراء
$\OLId{greek-variable}{Delta}(\sigma)\subseteq\OLId{greek-variable}{Delta}(\sigma')$ كلما تحقق
$\OLId{latin-looped}{R}\sigma\sigma'$.

\end{document}
